\chapter{职业伤害（职业安全）}
\setcounter{chapter}{9}
\chapterintro{
\textbf{【教学大纲要求】}

\imp{掌握：}职业安全的意义与任务；工伤流行病学；常见工伤事故及其危险因素；工伤事故的调查与评估；工伤事故预防对策与控制措施；职业安全的管理与事故预防对策。
}

\section{概述}

\begin{defbox}
\textbf{职业伤害（Occupational Injury）}：在生产劳动过程中，由于外部因素直接作用而引起机体组织的突发性意外损伤。

\textbf{职业安全}：以防止职工在职业活动过程中发生各种伤害和职业病为目的的工作领域及在法律、技术、设备、组织制度和教育等方面所采取的相应措施。
\end{defbox}

\section{安全事故的流行病学特征}

\begin{itemize}[leftmargin=*]
    \item \imp{事故发生率高的行业}：建筑、采矿、交通运输、制造业
    \item \imp{常见事故类型}：高处坠落、物体打击、机械伤害、触电、车辆伤害
    \item 时间分布：工作日的开始和结束时段事故高发
    \item 人群分布：青年工人、新工人、缺乏经验者事故率高
\end{itemize}

\section{导致工伤的因素}

\begin{enumerate}[leftmargin=*]
    \item 生产设备有缺陷
    \item 防护设备缺乏
    \item 劳动组织不合理、生产管理不善
    \item 个人因素（疲劳、注意力不集中、违反操作规程）
    \item 不良生产环境（照明不足、噪声等）
    \item 管理者与劳动者缺乏安全观念和知识
\end{enumerate}

\section{事故预防对策与措施}

\begin{enumerate}[leftmargin=*]
    \item 安全技术措施（设备安全设计、防护装置）
    \item 安全管理制度（操作规程、安全责任制）
    \item 安全教育与培训（三级安全教育）
    \item 个人防护用品的使用
    \item 事故应急处理预案
    \item 监督检查与事故调查分析
\end{enumerate}

\newpage

% =====================================================================
%  CHAPTER 10: 职业有害因素的识别与评价
% =====================================================================
\newpage